% vim: textwidth=78 ai spell spelllang=cs sw=2
\documentclass[11pt,a4paper]{article}

\usepackage[czech]{babel}
\usepackage{bibentry}
\usepackage[table]{xcolor}
\usepackage[utf8x]{inputenc}
\usepackage[IL2]{fontenc}
\usepackage[final]{hyperref}
\usepackage[all]{hypcap}
\usepackage{xspace}

%\setlength{\marginparwidth}{0pt}
% part of marginparwidth which is 106pt
%\addtolength{\oddsidemargin}{42pt}
%\addtolength{\evensidemargin}{-38pt}

\newcommand{\cdb}{\textsc{ClabureDB}\xspace}
\newcommand{\clang}{\textsc{Clang}\xspace}
\newcommand{\cppcheck}{\textsc{Cppcheck}\xspace}
\newcommand{\fbugs}{\textsc{FindBugs}\xspace}
\newcommand{\gcc}{\textsc{gcc}\xspace}
\newcommand{\klee}{\textsc{Klee}\xspace}
\newcommand{\ucklee}{\textsc{UC-Klee}\xspace}
\newcommand{\llvm}{\textsc{Llvm}\xspace}
\newcommand{\sage}{\textsc{Sage}\xspace}
\newcommand{\smatch}{\textsc{Smatch}\xspace}
\newcommand{\sparse}{\textsc{Sparse}\xspace}
\newcommand{\stan}{\textsc{Stanse}\xspace}
\newcommand{\symbio}{\textsc{Symbiotic}\xspace}
\newcommand{\perl}{\textsc{Perl}\xspace}
\newcommand{\xgcc}{\textsc{xgcc}\xspace}

\nobibliography*

\begin{document}
\pagestyle{empty}

\section*{Abstrakt}
Cílem této práce je výzkum aplikovatelnosti současných technik pro automatické
hledání chyb na velkých softwarových projektech. Vzhledem k tomu, že tyto
projekty bývají komplexní, je poptávka po výkonných analyzátorech, které mohou
takové projekty zpracovat bez přílišného generování falešných hlášení. Ke
splnění tohoto cíle se práce zaměřuje na statické analyzátory, zejména pak
na ty, které používají techniky jako abstraktní interpretace a symbolická
exekuce.  Od doby jejich uvedení v 70.\ letech 19.\ století se jejich použití
v současnosti zvyšuje díky rychlým procesorům. Jedině díky tomu jsou tyto
techniky použitelné na velkých softwarových projektech.

Práce shrnuje aktuální stav vybraných technik implementovaných v současných
statických analyzátorech. Nadále diskutuje několik velkých softwarových
projektů, které mohou být použity jako základna pro vyzkoušení technik.
Nakonec jako základnu pro zkoušení zvolíme Linuxové jádro. Kromě prozkoumaného
stavu současných technik navrhujeme v práci naše vlastní techniky. Práce
prezentuje nástroj \stan, který používá abstraktní interpretaci a který nemá
obtíže se zpracováním Linuxového jádra.  Nástroj také filtruje falešná hlášení
z některých inspektorů kódu. Filtrace je založena na vyhledávání vzorů, což se
ukázalo být nedostatečné. Proto jsme se následně snažili situaci zlepšit
symbolickou exekucí. K tomu účelu používáme symbolický exekutor \klee v našem
novém nástroji \symbio. Ale protože symbolická exekuce potřebuje dlouhý
výpočetní čas, nejprve kód prořezáváme technikou ,,static program slicing''.

Během našich experimentů jsme si vytvořili povědomí o kategoriích chyb, které
jsou obvykle přítomné v kódu jako je Linuxové jádro. Prozkoumali jsme je dále
a vyvinuli databázi známých chyb v jedné verzi jádra. Databáze je veřejně
dostupná a obsahuje informace o chybách a falešných hlášeních nalezených
několika nástroji hledajícími chyby. To může pomoci vývojářům nových
statických analyzátorů, aby si porovnali jejich zcela nový analyzátor oproti
databázi a zjistili, jak se jejich nástroj chová vzhledem k falešným hlášením,
k chybám, které minul a samozřejmě k chybám, které správně našel.

Zúčastnili jsme se také s nástrojem \symbio soutěže \emph{Software
Verification Competition}. Nástroj musel být upraven na potřeby soutěže. A
protože jsme se zúčastnili s příliš předčasnou verzí, výsledky nebyly příliš
povzbudivé. Nicméně, nástroj byl vylepšen a nyní dává dobré výsledky --
získáváme téměř všechny body v některých kategoriích soutěže. Plánujeme se
zúčastnit s naším nástrojem dalšího ročníku soutěže.

\clearpage

\section*{Autorův příspěvek}
\subsection*{Seznam publikací}
Na tomto místě uvádíme seznam publikací, ke kterým autor této práce přispěl
během svého Ph.D.\ studia. Snažíme se také uvést autorův podíl v~procentech.
Stanovit tento podíl není jednoduché, protože musíme uvážit několik kritérií,
např.\ kdo přišel s myšlenkou, kdo ji implementoval, kdo napsal publikaci,
prezentoval ji atd. Proto jsou procenta jenom subjektivní odhad.

Také prohlašujeme, že je tato práce založena na publikacích ~\ref{pub_svcomp},
\ref{pub_cdb}, \ref{pub_fmics} a \ref{pub_stanse} níže. Všechny byly
prezentovány autorem této práce na odpovídajících konferencích nebo
workshopech. Všechny z nich byly mezinárodní a uvedené články prošly
oponentským procesem. Tři z konferencí/workshopů jsou v~ERA označkování
uvedeny jako \emph{A}, \emph{B}, nebo \emph{C}. Další dvě publikace autora, na
nichž není tato práce postavena, byly publikovány na konferencích
označkováných \emph{A}.  Navíc, jedna z nich byla prezentována také autorem
této práce. Autor byl také \emph{pozván} a ochotně pozvání přijal, aby hovořil
o práci z publikace~\ref{pub_stanse} na konferenci ,,embedded world conference
2010''.  Autor byl také pozván, aby hovořil o vývoji Linuxového jádra na
workshopu NFX, který byl spoluorganizován Univerzitou Tomáše Bati ve Zlíně.

\begin{enumerate}
\item \bibentry{cse_atva}. [Podíl: 25\,\%]

\item \label{pub_svcomp} \bibentry{symbiotic-svcomp}. [Podíl: 75\,\%]

\item \label{pub_cdb} \bibentry{claburedb}. [Podíl: 80\,\%]

\item \label{pub_fmics} \bibentry{symbiotic}. [Podíl: 70\,\%]

\item \label{pub_stanse} \bibentry{stanse}. [Podíl: 40\,\%]

\item \bibentry{web2_isd}. [Podíl: 30\,\%]

\item \bibentry{web2_synasc}. [Podíl: 30\,\%]

\end{enumerate}

\subsection*{Seznam Softwaru}
Tady uvádíme přínos autora v softwarových projektech. Všechny projekty jsou
řádně popsány v této práci.

\begin{enumerate}
\item \stan
\item \symbio\ -- vyzdvihněme \llvm Slicer
\item \cdb
\end{enumerate}

\clearpage

\bibliography{cs}
\bibliographystyle{alpha}

\end{document}
